\documentclass[12pt]{report}
\usepackage[a4paper,margin=2.5cm]{geometry}
\usepackage{booktabs,tabularx,array,ragged2e,url,xltabular}
\newcommand{\sect}[1]{4.6.3}
\begin{document}
\noindent Filler.\par
\begingroup
\small
\renewcommand{\arraystretch}{1.25}
\begin{xltabular}{\textwidth}{@{}
    >{\RaggedRight\arraybackslash\hsize=0.62\hsize}X
    >{\RaggedRight\arraybackslash\hsize=1.38\hsize}X @{}}
\caption[Reproducibility record]{Reproducibility record: what is fixed, what is archived
and what is not recoverable.}\label{tab:Reproducibility} \\
\toprule
Item & Record \\
\midrule
\endfirsthead
\multicolumn{2}{@{}l}{\itshape Table \ref{tab:Reproducibility} continued} \\
\toprule
Item & Record \\
\midrule
\endhead
\midrule
\multicolumn{2}{r@{}}{\itshape continued overleaf} \\
\endfoot
\bottomrule
\endlastfoot
Code repository
  & Acquisition, screening, simulation, training and analysis in one repository, available
    at \url{https://github.com/Eric56123/Building_Sensor} and submitted alongside this
    dissertation \\
Language and libraries
  & Python 3.13.5; NumPy 2.4.2, SciPy 1.17.2, PyTorch 2.11.0, Matplotlib 3.10.8 \\
Training hardware & MacBook Pro (M2, 2022), 8\,GB memory \\
Deployment hardware
  & Raspberry Pi 4 Model B with ADXL345, as Table~\ref{tab:Reproducibility} \\
\addlinespace
Raw data
  & 311 ringdown captures across seven sessions, 0.14\,GB, archived with the submitted
    repository \\
Derived data
  & Simulated datasets regenerated from the forward model at run time; no intermediate
    array is required as an input \\
Regeneration
  & All classical modal results and all synthetic-network figures and tables regenerate
    from the raw captures through \texttt{make\_figures.py} with no manual step. The
    archived real-rig predictions of Section~\sect{} are the exception described below \\
\addlinespace
Seeds, data generation
  & Master seed 7 for the sampled-stiffness set, with a deterministic per-simulation seed
    derived from it; seed 0 for the retargeted training set \\
Seeds, training
  & 0 on weight initialisation in every run, so a weight ablation varies $\lambda$ alone.
    The seed sweep of Section~\sect{} is the exception, and varies the initialisation
    deliberately across twenty seeds \\
Solver reproducibility
  & The multi-start search and Monte Carlo perturbation of Section~\sect{} use a fixed
    seed; the start distribution, the $10^{-8}$~Hz branch tolerance and the perturbation
    scales are written alongside the results \\
Model record
  & The deployed retargeted model writes a record alongside its weights: fitted frequencies
    and stiffness ratios, training-set size, epoch count, loss weight and calibrated
    threshold \\
\addlinespace
Run provenance
  & Amplifier settings, sensor position, session date and joint start positions recorded
    per run, so each reported shift traces to the baseline it was measured against \\
Label audit
  & Plate, storey and folder name mapped in a single file, with an automated check of every
    damaged capture against its claimed location \\
\addlinespace
Known trap
  & \texttt{train\_rig\_pinn.py} writes its provenance record to the path given by
    \texttt{-{}-out}. Run without it, the previous record is overwritten rather than the run
    failing, and the sidecar is the artefact intended to survive \\
Known gap
  & The trained weights behind the rig predictions of Section~\sect{} were not retained.
    Those five predictions can be checked against the session record but not regenerated;
    retraining under the recorded configuration and seed reproduces the procedure, not
    necessarily that model \\
\end{xltabular}

\vspace{0.4em}
\noindent\begin{minipage}{\textwidth}
\footnotesize\RaggedRight
The classical modal results of Sections~\sect{} to~\sect{} regenerate in full from the
archived captures. The network results of Section~\sect{} regenerate on simulated data but
not, for the five rig cases, from the original weights.
\end{minipage}
\endgroup
\end{document}
